\IEEEPARstart{S}{tarting} from the breakthrough of AlexNet in large-scale image recognition~\cite{AlexNet}, deep learning has transformed classification, detection, and representation learning across scientific disciplines~\cite{DL-Survey}. Convolutional networks and vision transformers now provide complementary design choices to optimize performance for visual tasks~\cite{ViT,ConvNeXt}. To enhance effectiveness, developments ranging from BERT and generative pretraining in language~\cite{bert,gpt} to Swin transformers and YOLO in vision~\cite{swin,yolo} have utilized and depended on substantial training corpora. Datasets like ImageNet~\cite{imagenet}, Common Crawl~\cite{commoncrawl}, and OpenWebText~\cite{openwebtext} illustrate the breadth of this dependency. Acquiring a usable dataset is therefore a substantial part of the computational research process, rather than a negligible preliminary step.

The largest industrial laboratories can combine extensive storage, dedicated networking, and managed data-processing services. Academic groups more often assemble resources across institutional clusters, research clouds, and workstations. Even when training compute is available, data acquisition can consume a substantial portion of a project's preparation time. Provisioning more workers does not automatically solve this problem: data must be partitioned, transferred, checked, packaged, and made accessible to subsequent analysis. A practical system should support this workflow without assuming that every worker belongs to the same administrative environment.

Despite the breadth of accessible image repositories, converting URL-indexed collections into machine-learning-ready datasets remains demanding. iNaturalist, biodiversity portals, and curated web-image corpora distribute metadata that identifies externally hosted media assets. Users must handle various challenges such as changing availability, failed requests, duplicate rows, naming collisions, and storage layout. Metadata distribution is efficient for the publisher, but it transfers the work of acquisition and verification to each research group. The rate at which usable images become available directly affects the rate of research progress and model training.

Increasing request concurrency can improve utilization when individual transfers spend time waiting for remote responses. However, increasing it indiscriminately can also produce longer delays, HTTP refusals, and repeated work. Different public host services all differ in their capacity, policies, caching, and background demand. An HTTP error or delay increase is observable; the hidden mechanism that caused it usually is not. A useful acquisition tool should therefore treat these signals conservatively, respect explicit retry instructions, and maintain a record of what was actually delivered. It should not equate a high request count with a high rate of useful output.

Existing tools such as img2dataset provide effective URL-to-dataset pipelines, including parallel transfer, several output formats, configurable retries, and distributed execution~\cite{img2dataset}. However, these tools rely on configured process and thread counts rather than feedback-driven, per-host request admission, leaving concurrency tuning to the operator as source conditions change. FLOW-DC addresses these gaps by integrating per-host adaptive request admission, row-level outcome accounting, and artifact verification within an acquisition workflow designed for heterogeneous download resources.

\subsection{Automatic Concurrency Discovery under Unknown Server Conditions}
For a new source, an operator may know neither a useful concurrency setting nor how long that setting will remain appropriate. The best setting for one service or network path need not transfer to another. A conservative choice may leave capacity unused, whereas an aggressive choice can increase delay or provoke refusal. Repeating a manual search is both time-consuming and inconvenient, particularly when datasets span many hosts or when the available resources change between runs.

FLOW-DC makes this operating-point discovery part of acquisition itself. It begins with a small number of outstanding requests, observes application delay and overload feedback, and adjusts subsequent admission. This design can respond adaptively within a download job and can discover a new operating region when a previously useful baseline becomes stale. The practical value in this is that it reduces dependence on prior host knowledge and repeated operator intervention.

We take inspiration from Google's Bottleneck Bandwidth and Round-trip propagation time (BBR) congestion control~\cite{bbr}. BBR separates capacity probing from maintaining a low-delay reference. FLOW-DC carries these principles to HTTP acquisition, where the controllable quantity is the number of admitted requests rather than packet pacing. We introduce Policy-Aware Adaptive Request Control (PAARC), which combines a baseline-relative delay budget with the direction and persistence of excess-delay change to regulate request admission. PAARC requires neither access to server internals nor modifications to the transport stack.

FLOW-DC follows a manager--worker paradigm through TaskVine~\cite{taskvine}. A central manager partitions a dataset manifest and assigns acquisition tasks to workers. Each worker downloads asynchronously, applies the PAARC congestion rule, records row-level outcomes, and prepares verifiable output artifacts. Workers are easily able to join an ongoing workflow when their network access is compatible. This permits collaboration across resource pools without requiring a common filesystem or simultaneous startup.

The principal contributions are:
\begin{itemize}
\item An acquisition architecture that combines manifest partitioning, elastic task execution, asynchronous HTTP transfer, and verifiable output in one workflow.
\item An application-layer adaptive request controller, \paarc{}, that combines bounded exploration, a normalized excess-delay rate, and overload safeguards to discover useful concurrency settings.
\item A practical evaluation centered on verified goodput, completion time, and original-row coverage, together with applications in biodiversity and agricultural machine learning.
\end{itemize}
Section~\ref{sec:background} introduces the data and networking concepts. Section~\ref{sec:architecture} describes the architecture and controllers. Section~\ref{sec:results} presents the evaluation and discussion, and Section~\ref{sec:applications} describes research applications.
